%!TeX root = main.tex
The objective of this laboratory session is to parallelise the computation of the
covariance matrix using the Message Passing Interface (\texttt{MPI}) library. Unlike
the shared-memory model of \texttt{OpenMP}, \texttt{MPI} distributes work across
independent processes that communicate by explicitly passing messages, making it
suitable for both multi-core nodes and multi-node clusters. The results are compared
against the \texttt{OpenMP} implementation from the previous session, whose execution
times are reproduced in \Cref{tab:Ex_times_OpenMP} for reference.

\begin{table}[ht!]
    \centering
    \begin{tabularx}{\textwidth}{
        | >{\centering\arraybackslash}X
        | >{\centering\arraybackslash}X
        | >{\centering\arraybackslash}X
        | >{\centering\arraybackslash}X |
    }
        \hline
        \textbf{Number of Processes} & \textbf{Derived Variables} [s] & \textbf{Mean Values} [s] & \textbf{Cov.\ Matrix} [s]\\
        \hline
        1  & 62.82 & 2.26 & 16.12 \\
        2  & 31.64 & 1.08 &  9.42 \\
        4  & 12.82 & 0.72 &  6.51 \\
        8  &  6.74 & 0.53 &  5.34 \\
        16 &  3.57 & 0.48 &  3.33 \\
        20 &  3.84 & 0.51 &  3.13 \\
        30 &  2.80 & 0.46 &  3.60 \\
        40 &  2.65 & 0.52 &  3.59 \\
        \hline
    \end{tabularx}
    \caption{Execution times for different numbers of processes using \texttt{OpenMP}.}
    \label{tab:Ex_times_OpenMP}
\end{table}

\Cref{fig:Cov_OpenMP_graph} shows the covariance matrix execution time as a
function of thread count. Speedup is most pronounced up to 16 processes, after which
the curve flattens, indicating that synchronisation overhead and memory-bandwidth
saturation limit further gains.

\begin{figure}[ht!]
    \centering
    \begin{tikzpicture}
    \pgfplotsset{
        set layers,
        every axis legend/.style={
            cells = {anchor = east},% Centered entries
            anchor = north east,
            font = \footnotesize,
            shape = rectangle,
            fill = white,
            draw = black,
            at = {(1,1)}
            }
    }
    \begin{axis}[
        xlabel = {Number of processes[-]},
        ylabel = {Time [s]},
        ymin = 0,
        ymax = 65,
        xmin = 1,
        xmax = 40,
        symbolic x coords = {1,2,4,8,16,20,30,40},
        xtick = data,
        grid = both
        ]
        \addplot[red, mark=x, mark size=4pt] table[x=N_Threads, y=Cov_Matrix] {Data/openMP_CovMatrix.dat};
        \label{Covmat}
    \end{axis}
\end{tikzpicture}
    \caption{Covariance matrix execution time as a function of the number of
             \texttt{OpenMP} processes.}
    \label{fig:Cov_OpenMP_graph}
\end{figure}

\filbreak

\section{MPI Implementation and Performance}

The execution times obtained with the \texttt{MPI} implementation are reported in
\Cref{tab:Ex_times_MPI}. Compared to \texttt{OpenMP}, \texttt{MPI} achieves
consistently lower times for the derived-variables stage at high process counts,
reflecting the benefit of distributed memory parallelism. All processes read their
portion of the input file simultaneously using collective I/O routines
(\texttt{MPI\_File\_read\_at\_all}), which are designed to reduce I/O contention
on parallel filesystems. The execution times reported in \Cref{tab:Ex_times_MPI}
correspond to the \texttt{MPI~I/O} implementation. However, an anomalous increase
in the derived-variables time is observed at 80~processes (5.30\,s versus 1.30\,s
at 40), which is discussed in \Cref{sec:correctness}.

\begin{table}[ht!]
    \centering
    \begin{tabularx}{\textwidth}{
        | >{\centering\arraybackslash}X
        | >{\centering\arraybackslash}X
        | >{\centering\arraybackslash}X
        | >{\centering\arraybackslash}X |
    }
        \hline
        \textbf{Number of Processes} & \textbf{Derived Variables} [s] & \textbf{Mean Values} [s] & \textbf{Cov. Matrix} [s]\\
        \hline
        1  & 47.71 & 1.94 & 14.10 \\
        2  & 22.55 & 1.00 & 8.01  \\
        4  & 12.50 & 0.81 & 6.46  \\
        8  & 6.34  & 0.59 & 6.03  \\
        16 & 3.21  & 0.43 & 4.03  \\
        20 & 2.18  & 0.40 & 3.96  \\
        30 & 1.72  & 0.28 & 2.67  \\
        40 & 1.30  & 0.21 & 1.96  \\
        80 & 5.30  & 0.27 & 2.01  \\
        \hline
    \end{tabularx}
    \caption{Execution times for different numbers of processes using \texttt{MPI}.}
    \label{tab:Ex_times_MPI}
\end{table}

\begin{figure}[ht!]
    \centering
    \begin{tikzpicture}
    \pgfplotsset{
        set layers,
        every axis legend/.style={
            cells = {anchor = east},% Centered entries
            anchor = north east,
            font = \footnotesize,
            shape = rectangle,
            fill = white,
            draw = black,
            at = {(1,1)}
            }
    }
    \begin{axis}[
        xlabel = {Number of processes [-]},
        ylabel = {Time [s]},
        ymin = 0,
        xmin = 1,
        xmax = 80,
        symbolic x coords = {1,2,4,8,16,20,30,40,80},
        xtick = data,
        grid = both
        ]
        \addplot[red, mark=x, mark size=4pt] table[x=N_Threads, y=Read_Time] {Data/MPI_io_CovMatrix.dat};
        \label{read_MPI}

        \addlegendimage{/pgfplots/refstyle=read_MPI}\addlegendentry{MPI I/O}
    \end{axis}
\end{tikzpicture}
    \caption{Read times using \texttt{MPI}.}
    \label{fig:Read_MPI_graph}
\end{figure}

A notable anomaly appears in the read times of the \texttt{MPI I/O} implementation,
as shown in \Cref{fig:Read_MPI_graph}. At low process counts (1, 2, 4, and~8),
the collective I/O reads are significantly slower than the serial read performed
by rank~0 in the standard \texttt{MPI} implementation. This is likely caused by
the overhead of coordinating collective operations when the number of processes is
small relative to the file size, as well as filesystem-level locking when multiple
processes access contiguous blocks. Beyond 16~processes, the \texttt{MPI I/O}
read times decrease sharply and become competitive with the standard variant,
suggesting that the parallelism benefit only materialises once the per-process
I/O chunk is small enough to be read concurrently without contention.

\Cref{fig:Cov_MPI_graph} plots the covariance matrix execution time against the
number of \texttt{MPI} processes. A direct comparison with the \texttt{OpenMP}
results is given in \Cref{fig:Cov_comp_graph}.

\begin{figure}[ht!]
    \centering
    \begin{tikzpicture}
    \pgfplotsset{
        set layers,
        every axis legend/.style={
            cells = {anchor = east},% Centered entries
            anchor = north east,
            font = \footnotesize,
            shape = rectangle,
            fill = white,
            draw = black,
            at = {(1,1)}
            }
    }
    \begin{axis}[
        xlabel = {Number of processes [-]},
        ylabel = {Time [s]},
        ymin = 0,
        ymax = 65,
        xmin = 1,
        xmax = 80,
        symbolic x coords = {1,2,4,8,16,20,30,40,80},
        xtick = data,
        grid = both
        ]
        \addplot[red, mark=x, mark size=4pt] table[x=N_Threads, y=Remaining_Values] {Data/MPI_io_CovMatrix.dat};
        \label{Remval}
        \addplot[blue, mark=x, mark size=4pt] table[x=N_Threads, y=Mean_Values] {Data/MPI_io_CovMatrix.dat};
        \label{Meanval}
        \addplot[green, mark=x, mark size=4pt] table[x=N_Threads, y=Cov_Matrix] {Data/MPI_io_CovMatrix.dat};
        \label{Covmat}

        \addlegendimage{/pgfplots/refstyle=Remval}\addlegendentry{Derived Values}
        \addlegendimage{/pgfplots/refstyle=Meanval}\addlegendentry{Mean Values}
        \addlegendimage{/pgfplots/refstyle=Covmat}\addlegendentry{Covariance Matrix}
    \end{axis}
\end{tikzpicture}
    \caption{Covariance matrix execution time as a function of the number of
             \texttt{MPI} processes.}
    \label{fig:Cov_MPI_graph}
\end{figure}

In \Cref{fig:Cov_comp_graph}, the covariance matrix execution times for
\texttt{OpenMP} and \texttt{MPI} are plotted together for direct comparison.
Both implementations show a similar trend of decreasing execution time with
increasing thread/process count, but the \texttt{MPI} implementation consistently
outperforms \texttt{OpenMP} at higher counts, likely due to better scalability on
distributed memory systems. The ideal scaling lines based on the single-thread/process
execution times are also included for reference, illustrating the deviation from
perfect linear speedup as overheads become significant at higher parallelism levels.

\begin{figure}[ht!]
    \centering
    \begin{tikzpicture}
    \pgfplotsset{
        set layers,
        every axis legend/.style={
            cells = {anchor = east},% Centered entries
            anchor = north east,
            font = \footnotesize,
            shape = rectangle,
            fill = white,
            draw = black,
            at = {(1,1)}
            }
    }
    \begin{axis}[
        xlabel = {Number of processes [-]},
        ylabel = {Time [s]},
        ymin = 0,
        ymax = 20,
        xmin = 1,
        xmax = 40,
        symbolic x coords = {1,2,4,8,16,20,30,40,80},
        xtick = data,
        grid = both
        ]
        \addplot[red, mark=x, mark size=4pt] table[x=N_Threads, y=Cov_Matrix] {Data/openMP_CovMatrix.dat};
        \label{covmat_openMP}        

        % MPI data (threads 1–40 only)
        \addplot[blue, mark=x, mark size=4pt] table[x=N_Threads, y=Cov_Matrix] {Data/MPI_CovMatrix.dat};
        \label{Covmat_MPI}

        \addplot[red, dashed, thick]
            coordinates {
                (1,  16.116751)
                (2,  8.058376)
                (4,  4.029188)
                (8,  2.014594)
                (16, 1.007297)
                (20, 0.805838)
                (30, 0.537225)
                (40, 0.402919)
            };
        \label{is_openMP}

        \addplot[blue, dashed, thick]
            coordinates {
                (1,  14.096701)
                (2,  7.048351)
                (4,  3.524175)
                (8,  1.762088)
                (16, 0.881044)
                (20, 0.704835)
                (30, 0.467223)
                (40, 0.351611)
            };
        \label{is_MPI}

        \addlegendimage{/pgfplots/refstyle=covmat_openMP}\addlegendentry{openMP}
        \addlegendimage{/pgfplots/refstyle=Covmat_MPI}\addlegendentry{MPI}
        \addlegendimage{/pgfplots/refstyle=is_openMP}\addlegendentry{scaled openMP}
        \addlegendimage{/pgfplots/refstyle=is_MPI}\addlegendentry{scaled MPI}
    \end{axis}
\end{tikzpicture}
    \caption{Comparison of covariance matrix execution times for \texttt{OpenMP}
             and \texttt{MPI}.}
    \label{fig:Cov_comp_graph}
\end{figure}

\section{Correctness and Scalability\\of the MPI Implementation}
\label{sec:correctness}

The numerical correctness of the \texttt{MPI~I/O} implementation was verified by
comparing the covariance coefficients produced at each process count against the
serial reference (1~process). The reported entries are identical across all tested
configurations (2, 4, 8, 16, 20, 30, 40, and~80 processes) to the precision
printed by the program, confirming that the parallelisation introduces no
numerical discrepancy.

Performance, however, does not scale monotonically. The covariance matrix
computation time decreases from 14.10~s at 1~process to 1.96~s at
40~processes, but rises slightly to 2.01~s at 80~processes. A sharper
regression appears in the derived-variables stage, where the time increases from
1.30~s at 40~processes to 5.30~s at 80. This behaviour can be consistent
with over-decomposition: once the per-process workload becomes small, the cost of
collective operations such as \texttt{MPI\_Reduce} and inter-process
synchronisation exceeds the computational saving from additional parallelism.